% ============================================================
\chapter{Th\'eor\`emes de Bonnet-Myers et Synge}
\label{ch:bonnet_myers}
% ============================================================

\section{Introduction et motivations}

Peut-on d\'eduire la forme globale d'un univers \`a partir de mesures locales de sa courbure~? C'est pr\'ecis\'ement ce que font les th\'eor\`emes de Bonnet--Myers et de Synge, deux joyaux de la g\'eom\'etrie riemannienne globale. Le th\'eor\`eme de Bonnet (1855), red\'ecouvert et g\'en\'eralis\'e par Myers en 1941, affirme qu'une vari\'et\'e compl\`ete dont la courbure de Ricci est strictement positive est n\'ecessairement compacte et de diam\`etre fini. Autrement dit, la courbure positive ``referme'' l'espace sur lui-m\^eme. Le th\'eor\`eme de Synge (1936) va plus loin~: une vari\'et\'e riemannienne compacte de dimension paire et de courbure sectionnelle strictement positive est simplement connexe. Ces r\'esultats illustrent de mani\`ere spectaculaire comment la courbure contr\^ole la topologie.

\begin{intuition}
L'id\'ee directrice est que la courbure positive \og{}rapproche\fg les g\'eod\'esiques,
emp\^echant la vari\'et\'e de s'\'etendre \`a l'infini. Plus pr\'ecis\'ement~:
\begin{itemize}
    \item La courbure de Ricci positive uniform\'ement born\'ee inf\'erieurement
          force un \textbf{diam\`etre fini}, donc la compacit\'e.
    \item La courbure sectionnelle strictement positive en dimension paire
          contraint l'\textbf{orientabilit\'e} et la \textbf{simple connexit\'e}.
\end{itemize}
\end{intuition}

\section{Rappels sur les champs de Jacobi et la courbure de Ricci}

\begin{definition}[Courbure de Ricci]
\index{courbure de Ricci}
Soit $(M, g)$ une vari\'et\'e riemannienne de dimension $n$. La \emph{courbure de Ricci}
est la trace du tenseur de courbure~:
\[
    \operatorname{Ric}(u, v) = \sum_{i=1}^{n} g(R(e_i, u)v, e_i),
\]
o\`u $(e_1, \ldots, e_n)$ est une base orthonorm\'ee de $T_pM$.
\end{definition}

\begin{remarque}
La courbure de Ricci $\operatorname{Ric}(v, v)$ pour $\norm{v} = 1$ est la somme
des courbures sectionnelles $K(\sigma_i)$ des plans engendr\'es par $v$ et $e_i$,
pour $i \neq j$ o\`u $e_j = v$. C'est donc une \og{}moyenne\fg de la courbure sectionnelle.
\end{remarque}

\begin{definition}[Champ de Jacobi]
\index{champ de Jacobi}
Soit $\gamma : [0, L] \to M$ une g\'eod\'esique. Un \emph{champ de Jacobi} le long
de $\gamma$ est un champ de vecteurs $J$ le long de $\gamma$ v\'erifiant
l'\'equation de Jacobi~:
\[
    \frac{D^2 J}{dt^2} + R(J, \dot{\gamma})\dot{\gamma} = 0.
\]
\end{definition}

\begin{proposition}[Indice d'un champ de Jacobi]
\label{prop:indice_jacobi}
L'\emph{indice} d'une g\'eod\'esique $\gamma : [0, L] \to M$ est la forme bilin\'eaire~:
\[
    I(V, W) = \int_0^L \left[ g\!\left(\frac{DV}{dt}, \frac{DW}{dt}\right)
    - g(R(V, \dot{\gamma})\dot{\gamma}, W) \right] dt,
\]
d\'efinie sur les champs de vecteurs le long de $\gamma$ s'annulant aux extr\'emit\'es.
\end{proposition}

\section{Le th\'eor\`eme de Bonnet-Myers}

\begin{theoreme}[Bonnet-Myers]
\label{thm:bonnet_myers}
\index{th\'eor\`eme de Bonnet-Myers}
Soit $(M, g)$ une vari\'et\'e riemannienne compl\`ete de dimension $n \geq 2$.
Si la courbure de Ricci v\'erifie~:
\[
    \operatorname{Ric}(v, v) \geq (n-1)\kappa > 0
    \quad \forall\, v \in TM,\; \norm{v} = 1,
\]
alors~:
\begin{enumerate}
    \item Le diam\`etre de $M$ v\'erifie $\operatorname{diam}(M) \leq \frac{\pi}{\sqrt{\kappa}}$.
    \item $M$ est compacte.
    \item Le groupe fondamental $\pi_1(M)$ est fini.
\end{enumerate}
\end{theoreme}

\begin{proof}
\textbf{\'Etape 1~: Borne sur le diam\`etre.}
Soient $p, q \in M$ et $\gamma : [0, L] \to M$ une g\'eod\'esique minimisante
de $p$ \`a $q$, param\'etr\'ee par longueur d'arc. Supposons par l'absurde que
$L > \pi / \sqrt{\kappa}$.

Choisissons une base orthonorm\'ee $(e_1, \ldots, e_{n-1}, \dot{\gamma}(0))$ de
$T_pM$ et consid\'erons les champs de vecteurs obtenus par transport parall\`ele~:
$E_i(t) = P_{0,t}(e_i)$ pour $i = 1, \ldots, n-1$. D\'efinissons les variations~:
\[
    V_i(t) = \sin\!\left(\frac{\pi t}{L}\right) E_i(t).
\]
Chaque $V_i$ s'annule en $t = 0$ et $t = L$. Calculons la forme d'indice~:
\[
    I(V_i, V_i) = \int_0^L \left[ \frac{\pi^2}{L^2} \cos^2\!\left(\frac{\pi t}{L}\right)
    - \sin^2\!\left(\frac{\pi t}{L}\right) K(E_i, \dot{\gamma}) \right] dt,
\]
o\`u $K(E_i, \dot{\gamma})$ d\'esigne la courbure sectionnelle du plan $(E_i, \dot{\gamma})$.
En sommant sur $i = 1, \ldots, n-1$~:
\[
    \sum_{i=1}^{n-1} I(V_i, V_i) = \int_0^L \sin^2\!\left(\frac{\pi t}{L}\right)
    \left[ \frac{(n-1)\pi^2}{L^2} - \operatorname{Ric}(\dot{\gamma}, \dot{\gamma}) \right] dt.
\]
Si $\operatorname{Ric}(\dot{\gamma}, \dot{\gamma}) \geq (n-1)\kappa$ et
$L > \pi/\sqrt{\kappa}$, alors $(n-1)\pi^2/L^2 < (n-1)\kappa$, donc la somme
est strictement n\'egative. Ainsi au moins un $I(V_i, V_i) < 0$, ce qui signifie
que $\gamma$ n'est pas minimisante --- contradiction.

\textbf{\'Etape 2~: Compacit\'e.}
Comme $M$ est compl\`ete et de diam\`etre fini, elle est born\'ee. Par le
th\'eor\`eme de Hopf-Rinow, les ferm\'es born\'es d'une vari\'et\'e compl\`ete sont
compacts, donc $M$ est compacte.

\textbf{\'Etape 3~: Finitude de $\pi_1(M)$.}
Le rev\^etement universel $\widetilde{M}$ h\'erite de la m\'etrique pullback, qui
satisfait la m\^eme borne de courbure de Ricci. Par les \'etapes pr\'ec\'edentes,
$\widetilde{M}$ est compacte. Le rev\^etement $\pi : \widetilde{M} \to M$ est
donc fini, ce qui implique $\abs{\pi_1(M)} < \infty$.
\end{proof}

\begin{attention}[title=\textbf{Hypoth\`ese de compl\'etude}]
La compl\'etude est essentielle~: le disque ouvert $\mathbb{D}^2$ muni d'une m\'etrique
\`a courbure positive n'est pas compact. De m\^eme, l'hypoth\`ese $\kappa > 0$
(borne inf\'erieure \emph{strictement positive}) est cruciale~: un paraboloide
a $\operatorname{Ric} \geq 0$ mais est non compact.
\end{attention}

\begin{corollaire}[Comparaison avec la sph\`ere]
Si $(M^n, g)$ est compl\`ete avec $\operatorname{Ric} \geq (n-1)$, alors
$\operatorname{diam}(M) \leq \pi = \operatorname{diam}(S^n)$.
\end{corollaire}

\begin{exemple}[Sph\`ere et espaces projectifs]
La sph\`ere $S^n$ de rayon $r$ a $\operatorname{Ric} = \frac{n-1}{r^2} g$,
donc $\kappa = 1/r^2$ et la borne de Bonnet-Myers donne $\operatorname{diam}(S^n(r))
\leq \pi r$. L'\'egalit\'e est atteinte~: le diam\`etre de $S^n(r)$ est
exactement $\pi r$.
\end{exemple}

\begin{exemple}[Groupes de Lie compacts]
Tout groupe de Lie compact semi-simple muni de la m\'etrique de Killing (chang\'ee
de signe) a $\operatorname{Ric} > 0$. Ainsi $SU(n)$, $SO(n)$, $Sp(n)$ sont compacts
et ont un groupe fondamental fini.
\end{exemple}

\section{Le th\'eor\`eme de Cheng sur le diam\`etre maximal}

\begin{theoreme}[Cheng]
\label{thm:cheng}
\index{th\'eor\`eme de Cheng}
Soit $(M^n, g)$ une vari\'et\'e riemannienne compl\`ete avec
$\operatorname{Ric} \geq (n-1)\kappa > 0$. Si $\operatorname{diam}(M) = \pi/\sqrt{\kappa}$,
alors $M$ est isom\'etrique \`a la sph\`ere $S^n(1/\sqrt{\kappa})$.
\end{theoreme}

\begin{remarque}
Ce th\'eor\`eme de rigidit\'e montre que la borne de Bonnet-Myers est optimale et que
seule la sph\`ere atteint cette borne.
\end{remarque}

\section{Le th\'eor\`eme de Synge}

\begin{theoreme}[Synge]
\label{thm:synge}
\index{th\'eor\`eme de Synge}
Soit $(M, g)$ une vari\'et\'e riemannienne compacte \`a courbure sectionnelle
strictement positive ($K > 0$). Alors~:
\begin{enumerate}
    \item Si $\dim M$ est paire et $M$ est orientable, alors $M$ est simplement connexe.
    \item Si $\dim M$ est impaire, alors $M$ est orientable.
\end{enumerate}
\end{theoreme}

\begin{proof}[D\'emonstration (cas pair)]
Supposons $\dim M = 2m$ et $M$ orientable. Par l'absurde, supposons $\pi_1(M) \neq \{e\}$.
Il existe alors une g\'eod\'esique ferm\'ee $\gamma : [0, L] \to M$ de longueur
minimale dans sa classe d'homotopie libre non triviale.

Le transport parall\`ele le long de $\gamma$ d\'efinit une isom\'etrie
$P : T_{\gamma(0)}M \to T_{\gamma(0)}M$ pr\'eservant $\dot{\gamma}(0)$.
La restriction $P|_{\dot{\gamma}(0)^\perp}$ est une isom\'etrie du sous-espace
de dimension $2m - 1$ (impaire). Comme $M$ est orientable, $P$ pr\'eserve
l'orientation, donc $\det(P) = 1$. La restriction \`a $\dot{\gamma}(0)^\perp$
a $\det = 1$ aussi. Une isom\'etrie d'un espace de dimension impaire avec
d\'eterminant $1$ admet un vecteur fixe $v$ avec $\norm{v} = 1$.

Le champ de vecteurs parall\`ele $V(t)$ le long de $\gamma$ tel que $V(0) = v$
satisfait $V(L) = P(v) = v = V(0)$, donc $V$ est un champ de vecteurs parall\`ele
p\'eriodique perpendiculaire \`a $\dot{\gamma}$. Calculons~:
\[
    I(V, V) = -\int_0^L K(V, \dot{\gamma}) \, dt < 0,
\]
car $K > 0$. Cela signifie que $\gamma$ peut \^etre raccourcie par une variation
dans la direction $V$, contredisant la minimalit\'e de $\gamma$ dans sa classe
d'homotopie.
\end{proof}

\begin{corollaire}
Les seules vari\'et\'es de dimension paire \`a courbure sectionnelle strictement
positive et groupe fondamental non trivial sont les espaces projectifs r\'eels
$\R P^{2m}$ (non orientables).
\end{corollaire}

\begin{exemple}[$S^{2m}$ et $\C P^m$]
La sph\`ere $S^{2m}$ ($K = 1$) est simplement connexe, en accord avec Synge.
L'espace projectif complexe $\C P^m$ ($K \in [1, 4]$) est orientable de dimension
paire $2m$, et effectivement $\pi_1(\C P^m) = 0$.
\end{exemple}

\begin{exemple}[$S^3$ et $\R P^3$]
En dimension impaire, Synge ne force pas la simple connexit\'e~: $\R P^3$ a $K > 0$
et $\pi_1(\R P^3) = \Z/2\Z \neq 0$, mais il est orientable, en accord avec le
th\'eor\`eme.
\end{exemple}

\section{Applications et cons\'equences topologiques}

\begin{proposition}[Contraintes sur le groupe fondamental]
Soit $(M^n, g)$ compl\`ete avec $\operatorname{Ric} \geq (n-1)\kappa > 0$. Alors~:
\begin{enumerate}
    \item $\pi_1(M)$ est fini (Bonnet-Myers).
    \item Si $n$ est pair et $M$ est orientable, $\pi_1(M) = 0$ (Synge).
    \item Si $n$ est impair et $K > 0$, $M$ est orientable (Synge).
\end{enumerate}
\end{proposition}

\begin{theoreme}[Borne de volume -- Bishop-Gromov]
Si $(M^n, g)$ est compl\`ete avec $\operatorname{Ric} \geq (n-1)\kappa$, alors pour
tout $p \in M$ et $0 < r \leq R$~:
\[
    \frac{\operatorname{Vol}(B(p, R))}{\operatorname{Vol}(B(p, r))}
    \leq \frac{V_\kappa(R)}{V_\kappa(r)},
\]
o\`u $V_\kappa(r)$ est le volume de la boule de rayon $r$ dans l'espace \`a courbure
constante $\kappa$.
\end{theoreme}

\begin{formules}[title=\textbf{R\'esum\'e des bornes de courbure et topologie}]
\begin{center}
\renewcommand{\arraystretch}{1.3}
\begin{tabular}{lll}
\textbf{Hypoth\`ese} & \textbf{Conclusion} & \textbf{Th\'eor\`eme} \\
\hline
$\operatorname{Ric} \geq (n-1)\kappa > 0$, compl\`ete
    & $\operatorname{diam} \leq \pi/\sqrt{\kappa}$, compacte
    & Bonnet-Myers \\
$\operatorname{Ric} \geq (n-1)\kappa > 0$, compl\`ete
    & $\pi_1(M)$ fini
    & Bonnet-Myers \\
$K > 0$, compacte, $\dim$ paire, orientable
    & $\pi_1(M) = 0$
    & Synge \\
$K > 0$, compacte, $\dim$ impaire
    & $M$ orientable
    & Synge \\
$\operatorname{Ric} \geq (n-1)\kappa$, compl\`ete
    & Comparaison de volume
    & Bishop-Gromov \\
\end{tabular}
\end{center}
\end{formules}

\section{Exercices}

\begin{exercice}
Montrer que $\R P^n$ muni de la m\'etrique quotient de $S^n$ a $K = 1$. Pour $n$
pair, v\'erifier que $\R P^n$ n'est pas orientable et a $\pi_1(\R P^n) = \Z/2\Z$,
en accord avec Synge.
\end{exercice}

\begin{exercice}
Soit $(M^n, g)$ compl\`ete avec $\operatorname{Ric} \geq (n-1)$. Montrer que
le volume de $M$ v\'erifie $\operatorname{Vol}(M) \leq \operatorname{Vol}(S^n)$.
\end{exercice}

\begin{exercice}
Soit $M = S^2 \times S^2$ munie de la m\'etrique produit. Calculer la courbure de
Ricci et v\'erifier que $\operatorname{Ric} > 0$. En d\'eduire que $S^2 \times S^2$
est compacte et simplement connexe (par Synge, $\dim = 4$ paire).
\end{exercice}

\begin{exercice}[Optimalit\'e de Bonnet-Myers]
Construire une famille de m\'etriques sur $S^n$ avec $\operatorname{Ric} \geq (n-1)\kappa$
pour un $\kappa$ donn\'e, et dont le diam\`etre tend vers $\pi/\sqrt{\kappa}$.
\end{exercice}

\begin{exercice}
Montrer que le tore $\mathbb{T}^n$ n'admet aucune m\'etrique \`a courbure de Ricci
strictement positive. \emph{Indication~:} utiliser $\pi_1(\mathbb{T}^n) = \Z^n$.
\end{exercice}

\begin{exercice}
Soit $(M, g)$ une vari\'et\'e riemannienne compacte de dimension $3$ avec $K > 0$.
Montrer que $M$ est orientable. Peut-on conclure que $M$ est simplement connexe~?
Justifier avec l'exemple de $\R P^3$.
\end{exercice}
